\section{Extended experimental analysis}
\label{app:revision-details}


We report the completed Codex runs in the inventory generated on 24 September 2026,
separating the broader GPT-5.6 Luna evaluation from the recent ablations. The inventory
counts completed attempts, including unsuccessful submissions and stopping-limit
terminations. For the ablations, all 360 planned sample--configuration pairs enter
the denominator. Execution errors and omissions after blocking errors count as
unsuccessful evaluations, with causes marked separately. We retain the first completed
attempt per sample and configuration when available; retries do not add denominator entries.
% Source: workspace/RESULTS/run-inventory/report.md (2026-09-24T23:13:33.275Z).
% Generated with workspace/report_run_inventory.mjs; per-attempt provenance is in report.json.
% The prior Results section and its editorial notes are preserved in
% sections/09-results-previous-snapshot.tex, but are not included in the paper.


Earlier GPT-5.6 Luna runs used different samples and budgets per condition, so we list them
without comparison (Table~\ref{tab:headline}); tool effects are measured only in the ablations.

\subsection{Tool and reasoning-effort ablations}



The recent study uses the first ten samples of each group in every configuration:
\texttt{easy-}, \texttt{mid-} and \texttt{hard-0001}--\texttt{0010} from the all-layers corpus and
\texttt{layers-0001}--\texttt{0010} from the some-layers corpus; the some-layers samples all have five reference folds. It compares GPT-5.6 Luna at low effort with
GPT-6 Luna at low and high effort. Tool conditions and run settings are
given in Sections~\ref{sec:ablations} and~\ref{sec:exp-protocol}. Of the 360 planned
sample--configuration pairs, five are protocol failures counted as unsolved
(Appendix~\ref{app:protocol-failures}).

\paragraph{Sample selection.} IDs follow consecutive generator seeds and were fixed before any
ablation run; the selected samples match their groups in depth but have somewhat fewer layers and
creases, so their solve rates if anything overstate the full groups' (Appendix~\ref{app:sample-selection}).



\paragraph{Tools.} Summing the Total column of Table~\ref{tab:luna-ablations} over the three model
settings gives, for the same 40 samples in each, 8/120 solves with B, 15/120 with L and 27/120
with L+A (two L and three L+A pairs are protocol failures counted as unsolved); the order B$<$L$<$L+A
holds in each setting. By group, easy gives 5, 2 and 8 of 30 and some-layers 3, 13 and 19 of 30;
medium and hard give none. These are descriptive comparisons over reused instances. In particular, the feedback rows
were generated under a defective comparison tool, whose trajectory-level effect has unknown
sign. Terminal re-scoring does not remove that confound; we do not interpret these totals as
a controlled estimate of the benefit of corrected feedback.

\paragraph{Feedback without enumeration.}

Because L+A adds two things to B, the ablation above cannot say whether target feedback helps
without legal enumeration. We therefore ran B+A (basic tools plus automatic tier-3 feedback, no
\texttt{list\_legal\_folds}) for GPT-5.6 Luna at low effort on the same 40 instances, with the
image ablation's settings (text and images, 80 turns, one episode per instance;
Table~\ref{tab:feedback-only}). All five B+A solves also satisfy $E_{\mathrm{id}}$ on replay.
Against B (2/40), feedback alone gives 5/40, the same as legal enumeration alone (5/40), and the
two together give 9/40. With ten instances per group none of these differences is significant on
its own. The totals do not decompose by instance: B+A solves easy-0001, 0004 and 0006, and at most
one of these is among the L+A easy solves; L solves layers-0007, which L+A does not. With one
episode per instance, which instances a condition solves varies between runs: five further repeats of the
text-only L+A run solved 2, 3, 1, 3 and 3 easy instances under $E_{\mathrm{id}}$, against 5 in
the original run, with different instance IDs (Section~\ref{sec:limitations}). We therefore read Table~\ref{tab:feedback-only} only as showing
that each addition raises the total, not that their effects add or which instances each helps.



\paragraph{Reasoning effort.} Pooled over the three tool conditions, GPT-6 Luna solves 6/30
some-layers samples at low effort and 18/30 at high effort. Per condition the increases are
0/10 to 3/10, 1/10 to 7/10 and 5/10 to 8/10; with ten samples per cell only the second is
significant on its own (two-sided Fisher exact test, $p\approx0.02$; $p\approx0.21$ and $0.35$
for the others), and a 95\% Wilson interval spans roughly $\pm0.3$ (e.g.\ 0.11--0.60 for
3/10). Pairing outcomes by sample and condition, every discordant pair favours high effort
(12 to 0; exact McNemar test, $p<0.001$), so the effect is consistent even where single cells
are underpowered. McNemar treats the 30 sample--condition pairs as independent, but each
some-layers instance appears in all three tool conditions. Treating the instance as the unit,
a cluster bootstrap over the ten instances gives a high-minus-low difference of
0.40 in some-layers solve rate (95\% CI 0.23--0.53), and an exact sign-flip test over
per-instance differences gives $p=0.008$. This instance-clustered result concerns the ten five-fold some-layers instances and the
observed historical tool conditions; it is not a clean estimate under corrected feedback.
The effect magnitude is uncertain and does not establish performance at larger depths.

\paragraph{Depth.} No Luna ablation solves a medium or hard sample, and the breadth-first
baseline solves 1/40 medium and 0/16 hard attempts (Appendix~\ref{app:bfs-analysis}), against
366/669 on the easy group (at most four of which fail $E_{\mathrm{id}}$; Appendix~\ref{app:identity-audit}). Within the easy group, the model solve rate falls from 32\% at four folds to
0--15\% at eight to ten (Table~\ref{tab:easy-depth}). The groups differ in more than depth. In
the all-layers tier every fold moves the whole stack and usually crosses every layer, so the
layer count grows roughly exponentially with depth: over the 40 ablation instances the final
stack has a median of 23 layers on easy (maximum 64), 90 on medium (448) and 512 on hard (3,504),
and over the full released groups the hard median is 404 (maximum 16,384). Depth and layer count
are therefore confounded, and the hard group mostly measures states of hundreds to thousands of
layers, whose serialised geometry approaches or exceeds what fits in a prompt (one medium
episode failed on input size; Appendix~\ref{app:protocol-failures}). These results show that the
deeper groups are harder, not that depth alone is the cause, and we do not read the hard group as
a limit of planning depth (Section~\ref{sec:limitations}).
Breadth-first search time on the saved searches grows by a fitted factor of 3.9 per fold on easy
and 4.6 on medium (Figure~\ref{fig:depth-wall}); the hard-group factor of 12.4 rests on two samples
and an extrapolation, and we do not use it as a measure of task complexity.

\paragraph{Matched search baselines.} The comparison above mixes budgets and instance sets. To
compare the model with search on the same footing, we ran two baselines that use no model on the
same 40 ablation instances, with the action set the model is given and success defined as CP match
plus $E_{\mathrm{id}}$, checked after every fold (Table~\ref{tab:matched-baselines}). \emph{Random}
plays 80 uniformly random legal folds per episode (the model's turn budget), 100 episodes per
instance. \emph{BFS} searches distinct states (identity-keyed) up to the reference depth plus four,
and we record whether it reaches the target within 80, $10^3$, $10^4$ and $10^5$ expansions (600\,s
wall clock per instance). Within an 80-fold budget, random legal folding solves 57\% of easy and 96\%
of some-layers episodes, while the best model conditions solve at most 5 of 10 easy instances
(GPT-5.6 Luna, text-only L+A) and 8 of 10 some-layers instances (GPT-6 Luna, high effort, L+A)
in one episode each (Table~\ref{tab:luna-ablations}, Section~\ref{sec:results-matched}).
BFS solves 5 of 10 medium instances within $10^5$ expansions; random folding solves at least one
episode on 3 of 10 medium instances. Only the hard group is unsolved by every method we ran, and
its states are the largest (above). We therefore do not describe medium as beyond reach: on easy,
some-layers and medium the model is less effective than undirected search over the same legal
actions. Search with $10^5$ expansions and 100 random episodes per instance uses far more
attempts than one model episode; the comparison shows that the model's choices carry less
information than chance would over the same number of folds, not that search is cheaper or
better in general.



\subsection{Input modality}



To measure what the images contribute, we reran the GPT-5.6 Luna (low effort) ablation with
every image removed from the prompt and from tool results, keeping the JSON geometry of the crease
pattern, target and current state. Under $E_{\mathrm{id}}$, text-only solves 3, 5 and 11 of 40 with
B, L and L+A, against 2, 5 and 9 with images (pooled 19/120 vs.\ 16/120); under the submitted
$E_{\mathrm{plane}}$ scoring the text-only counts were 3, 7 and 11 (21/120 vs.\ 18/120; discordant
pairs 9 vs.\ 6, exact McNemar $p=0.61$). Removing images does not reduce solve rate on these
samples. The prompt already states every vertex coordinate, parity and layer rank exactly, while a
512-pixel rendering shows the same geometry less precisely and cannot show which layer lies above
which; the model's successful reasoning runs through the coordinates. At this sample size we
cannot exclude a small benefit from images, but we find none. CP2Seq as configured here therefore
measures geometric reasoning over an exact numeric description of every layer (coordinates,
orientation, stacking order and sheet region) more than visual perception. One
text-only basic and two text-only medium episodes ended in protocol failures and count as unsolved.

An images-only condition, which keeps the renderings and removes the coordinates, isolates
perception (Appendix~\ref{app:images-only}). With basic tools it solves 1/40 (one easy instance),
against 3/40 text-only and 2/40 with both (Table~\ref{tab:modality}). The differences are one and
two instances out of 40 and are not significant; they are consistent with the text-only result in
that images do not substitute for coordinates.



\subsection{Failure analysis}


Termination labels name the stopping rule that fired, not the cause of failure. Rejected
actions never end an episode, although they consume turns; the cycling rule stops an episode
when the same fold is played from the same position three times. In the ablations, 191 of the 305
unsuccessful completed episodes ($E_{\mathrm{id}}$; 300 under $E_{\mathrm{plane}}$) end under the cycling rule (Table~\ref{tab:ablation-terminations}); in the
earlier 447-attempt GPT-5.6 Luna snapshot, 215 of 384 do. These counts describe how episodes
stopped, not why they failed, and we do not attribute failures to capabilities from them.

The per-turn logs of the image ablation episodes were not retained; the released runs keep each
episode's accepted sequence and final state. Per-turn logs survive for the text-only GPT-5.6 Luna
episodes, and we analysed those turn by turn (Table~\ref{tab:trajectories}), replaying each
accepted state and comparing it with the target. Each unsuccessful episode is assigned one of four
mechanical categories (Appendix~\ref{app:matched-trajectories}): F1, mostly geometrically illegal
proposals; F2, mostly legal folds off the target pattern; F3, the full pattern at the target depth
but a wrong stack; F4, never reaching the target depth. With basic tools, 82\% of proposed folds
are rejected, mostly as off-pattern, and 29 of 36 failures are F2; with legal enumeration, the
rejection rate falls to 13\% and failures shift to F4 (21 of 32) and F3 (6 of 32). These failure
counts use the outcomes as submitted ($E_{\mathrm{plane}}$), as in Table~\ref{tab:trajectories}. Once the tool
removes illegal and off-pattern moves, the remaining failures are episodes that never reach the
target depth, or reach it with the layers in the wrong order. The categories locate where an
episode broke down, not which internal capability failed, and they come from one model without
images; we have not verified that they transfer to the image ablations.

\paragraph{Reaching the target without stopping.} An episode could fail by passing through the
target state and folding on instead of finishing. Replaying every accepted state of the 143 saved
episodes with per-turn logs and checking CP match plus $E_{\mathrm{id}}$ after each, an automatic
stop at the first match would raise solves from 25 to 27. Failures are therefore almost never
missed stops; in all but two unsuccessful episodes the model never reaches the target.

\paragraph{Exploration.} We counted the distinct folded states (flat sheet excluded) each saved
episode reaches and compared them with 20 random episodes of 80 legal folds on the same instance.
The median per-instance ratio of model to random distinct states (easy, some-layers) is 0.11 and
0.12 with text-only B, 0.16 and 0.24 with text-only L, and 0.25 and 0.18 with text-only L+A; it is
lowest for images-only B (0.12 and 0.07). The model reaches between a tenth and a quarter of the
states random folding reaches, and more states with enumeration than without, as with solve rate.
With basic tools 72--81\% of proposals are rejected; with legal enumeration 57--65\% of accepted
folds return to a state already reached. Part of the gap reflects episode length: model episodes
stop after 20--45 turns on average (cycling rule or \texttt{finish}), whereas every random episode
uses all 80 folds. Together with the matched baselines, this indicates that the model explores
far less of the legal state space than undirected search and revisits what it has seen, rather
than failing at the final step.

\paragraph{Backtracking.} We examined every backtrack (a \texttt{go\_to\_step},
\texttt{restore\_revision} or \texttt{remove\_fold} that shortens the sequence) in the saved easy and
some-layers episodes, and for each asked whether a local repair existed: whether the target is
reachable (CP match plus $E_{\mathrm{id}}$) from the sequence with its last one, two or three folds
removed, by depth-limited search within the reference depth plus two (node cap 4,000). In the
120 text-only GPT-5.6 Luna L+A episodes (the original run and the five repeats) there are 384
backtracks; 104 (27\%) return to the flat sheet, and the mean backtrack removes 2.6 folds. The
search decided all 384: a repair within three folds existed for 232 (60\%). A repair existed for
71 of the 104 returns to the flat sheet (68\%) and for 161 of the 280 shorter backtracks (58\%), so
the model restarts about as often when a short repair is available as when it is not. In 21 Claude
Sonnet~5.5 L+A episodes, 16 of 82 backtracks (20\%) return to the flat sheet; a repair existed for
8 of these 16 (50\%) and for 36 of the other 66 (55\%).

After a backtrack of two or more folds, we asked whether the next accepted folds re-make the
discarded ones, reaching the same states in the same order, and compared this with a policy that
picks uniformly among the distinct legal folds at each state. With legal enumeration the models
re-make the discarded folds less often than chance: GPT-5.6 Luna L+A in 32 of 293 cases against
61 expected, and Claude Sonnet~5.5 in 0 of 58 against 8.7 expected. With basic tools, where no
list is shown, GPT-5.6 Luna re-makes them more often than chance (4 of 6 against 0.9 expected with
B; 8 of 14 against 3.7 with B+A), but these counts are small.

The two families differ in whether they return to branches they have already given up. Counting
every state on a discarded suffix, and every state reported to have no legal fold, as abandoned,
37\% of GPT-5.6 Luna L+A add\_folds after an episode's first backtrack land on an abandoned state
again (340 of 931; 65 of them on a state with no legal fold), against 11\% for Claude Sonnet~5.5
(26 of 243; 5 on such states); 77 of 120 Luna episodes and 7 of 21 Sonnet episodes contain such a
revisit. Each turn is a fresh request, so the only record of abandoned branches is the action
history in the prompt; Sonnet's lower rate is consistent with it using that history to avoid
branches already tried, but its reasoning text is not returned by the interface we used, so we
cannot tell how. Re-entering an abandoned shallow state is not always an error, since a solution
may pass through it; re-entering a state with no legal fold always is.

The tier-3 comparison can mark a state as close to the target when no short repair exists. On
\texttt{layers-0003}, GPT-5.6 Luna (nine restarts in six episodes) and Claude Sonnet~5.5 (one)
restarted from a state with 8 of 9 layers, 6 of the 8 matching the target; in seven of these ten
cases the search found no repair within three folds (in the other three a one-fold repair
existed), and in the Sonnet case the state had no legal fold. On \texttt{easy-0002}, five Luna
episodes and one Sonnet episode reached the target's 24 layers in the wrong order and restarted;
none had a repair within three folds. Restarting from such states is therefore not by itself an
error. What the episodes show is that the models do not distinguish a recoverable near miss from a
trap: whether a short repair exists barely changes how often they restart, while the feedback's closeness signal (layer count and matching layers)
also marks unrecoverable states as close. These counts come from a bounded search; they describe
behaviour, not its cause.

\paragraph{A second model family.} Claude Sonnet~5.5 ran the same text-only L+A configuration on
the 20 instances of the repeated-trial set, one episode each (the first four easy episodes used an
earlier prompt order). Under $E_{\mathrm{id}}$ it solves 5 of 10 easy and 9 of 10 some-layers
instances (14/20); under $E_{\mathrm{plane}}$ it solves 7 of 10 easy, and the two
$E_{\mathrm{plane}}$-only solves ended with \texttt{finish} after the comparison tool reported
\emph{matches}, as for GPT-5.6 Luna (Section~\ref{sec:limitations}). The unsolved some-layers
instance, \texttt{layers-0006}, is a protocol failure: in two attempts a single call exceeded the
300-second timeout (at turns 31 and 37); GPT-5.6 Luna never solved it in five repeats. On the same
20 instances GPT-5.6 Luna solves 8.4 per repeat on average (2.4 easy, 6.0 some-layers) and 14 in at
least one of five repeats. One Sonnet episode per instance thus matches Luna's five-repeat
coverage, but the solved easy instances differ (Sonnet: 0001, 0004, 0006, 0009, 0010; Luna in at
least one repeat: 0002, 0005, 0006, 0009, 0010). With one episode per instance and 20 instances
this is consistent with a stronger model, not a measured difference; Sonnet still falls short of
random legal folding on the same groups (Table~\ref{tab:matched-baselines}).





The broader evaluation, ablations, BFS snapshot and CP-distance scan cover different attempt
sets with different budgets; we report their denominators separately and do not pool them into
a single controlled comparison. The one matched comparison of models and search is
Table~\ref{tab:matched-baselines}, on the 40 ablation instances.

\section{Extended limitations and repeated-trial accounting}
\label{app:extended-limitations}

The main ablations cover Luna models and BFS, with an added 20-instance Sonnet evaluation, and the ablations use ten samples per group and
configuration. This study is currently privately funded, and these limits follow from the cost of
evaluation: with full image history, one Luna episode consumed about 3.1 million input tokens on
average over 20 episodes. In pilots, Qwen3.5-9B (via Tinker) cycled without solving either of two samples, and one
80-turn GPT-6 Astra episode cost about US\$400. We are seeking research funding to extend the evaluation to further
model families and to the full corpus.

The action space covers all-layers and selected some-layers simple folds, excluding
pre-creasing; the some-layers corpus currently extends to six reference folds. The environment
represents zero-thickness flat states and does not model material mechanics or continuous
collisions during folding. CP2Seq is a planning benchmark over an exact state space, not a model
of physical folding, and we draw no engineering conclusions from it.

\paragraph{Depth and layer count are confounded in the all-layers tier.} Whole-stack folds sampled
uniformly almost always cross every layer, so the released medium and hard groups have medians of
94 and 404 final layers (maximum 16,384), and 96 of the 100 hard samples end with a footprint under
5\% of the sheet. Results on these groups, including every model's and search's failure on hard,
cannot be attributed to sequence depth alone. The generator already provides a cap on the number
of layers a single fold may cross; with a cap of 16, all three groups fill their quotas with medians
of 24, 60 and 76 final layers (maximum 120). The results in this paper are on the uncapped release;
re-evaluation on the capped corpus is future work.

\paragraph{One episode per instance.} The ablations run one episode per instance and configuration.
Five further repeats of the text-only GPT-5.6 Luna L+A run solved 2, 3, 1, 3 and 3 easy instances
under $E_{\mathrm{id}}$, against 5 in the original run, with different instances solved. Both the totals and the set of solved instances therefore carry
run-to-run variance that one episode cannot estimate; we interpret only differences in totals, with
their uncertainty, and not per-instance differences between configurations. Repeated episodes per
instance were then run for the text-only GPT-5.6 Luna L+A configuration: five episodes on each
of the 10 easy and 10 some-layers instances (100 episodes). Under $E_{\mathrm{id}}$ the five
repeats solve 7, 10, 6, 10 and 9 of 20 (42/100 episodes), and 14 of 20 instances are solved in at
least one repeat (easy 5/10, some-layers 9/10); six instances are never solved. Under
$E_{\mathrm{plane}}$ the counts are 47/100 episodes and 16/20 instances (easy 7/10). The five
episodes counted as solved only under $E_{\mathrm{plane}}$ (\texttt{easy-0004} in four repeats,
\texttt{easy-0003} in one) each ended with the model calling \texttt{finish} after the comparison
it was shown reported \emph{matches} with every layer agreeing, so they record the plane-only
comparison tool, not a model error. The gap between per-episode rate and coverage measures
run-to-run variance; it is obtained with five times the episode budget and is not evidence of more
efficient search.

We treat these repeated trials as a stability analysis, not a causal diagnosis of failure.
A planned diagnostic study freezes the task, state, history and reasoning effort, then varies
one form of support: a summary of previously observed attempts, explicit action semantics, or
uniform one-step previews. Its controls share the same legal-action menu. Such interventions
can test the effect of support on decisions and subsequent success; they do not by themselves
identify an internal memory or reasoning deficit. This study has not yet been run.

\paragraph{Input modality.} Removing rendered images does not lower solve rate, and an images-only
condition solves 1/40 with basic tools. The benchmark as configured measures reasoning over exact
coordinates; it is not a test of visual perception.

Termination labels record stopping conditions. The trajectory categories of
Section~\ref{sec:results-failures} locate where an episode broke down, but they are mechanical
rules over observable states and do not identify which internal capability failed. Revisit and
budget limits constrain the observed search, and the effect of a looser cycling limit is untested.
Per-turn logs of the image ablations were not retained, so the trajectory analysis covers the
text-only runs only, and the comparison-tool defect cannot be audited call by call there.
Training overlap with the public corpus is unknown.

Corpus replay checks assess consistency with the shared generation engine. The submitted results
were scored with a plane-only terminal comparator that did not check original-sheet region
identity; the identity audit (Appendix~\ref{app:identity-audit}) found accepted solutions that
fail the identity-aware relation, and all results are now re-scored under it. The comparison
tool shown to models still uses the plane-only relation, so its \emph{matches} can be wrong in
the same way; the affected calls are counted in Appendix~\ref{app:protocol-failures}.
